\section{Design of SCOUT}
\label{sec:design}


\subsection{Probe Structure and Diagnostic Pipeline}
\label{sec:overview}

\paragraph{Probe structure}
SCOUT uses the same fixed-size probe structure at both
replicas.  Each replica maintains $V$ independent \emph{views};
each view is a linear array of $L$ \emph{sections}, each
section containing $m_0$ cells. Each cell stores a 32-bit
signed counter, so each view contains $M = L \cdot m_0$ cells.
Across all $V$ views, the probe occupies $4VLm_0$ bytes
(e.g., ${\sim}512$\,KiB at $V{=}2$, $L{=}2048$, $m_0{=}32$),
fixed at configuration time and independent of data scale~$P$.

A two-level mapping (Fig.~\ref{fig:scout-mapping-structure-a})
writes each key into this array: the key's segment determines
a contiguous window of $W$ sections on the section axis, and a
second hash selects $k$ cells within that window.  The
contiguous-window structure is central to candidate
compression and is analyzed in
\S\ref{sec:spatial-coupling}.
For analysis, both mapping stages are modeled as uniform and
independent hashing; hash robustness is verified in Supplement~\ref{sec:appendix-hash}.

\begin{figure}[!htbp]
\centering
% SCOUT probe structure (a): two-level mapping (No in-figure title)
\begin{tikzpicture}[
    x=1cm, y=1cm, font=\footnotesize, >=Stealth,
    seg/.style    = {draw, rounded corners=1pt, fill=black!5,
                     inner sep=2pt, minimum height=0.40cm},
    segHi/.style  = {seg, draw=blue!65!black, line width=0.8pt, fill=blue!8},
    segSmall/.style = {draw, rounded corners=1pt, fill=black!5,
                       inner sep=0pt, minimum width=0.44cm,
                       minimum height=0.32cm, font=\scriptsize},
    cell/.style   = {draw, inner sep=0pt, minimum size=0.26cm},
    lbl/.style    = {font=\scriptsize, align=center, inner sep=1.5pt},
    note/.style   = {font=\scriptsize, text=black!60},
    arrlbl/.style = {font=\scriptsize, align=center, fill=white,
                     fill opacity=.92, text opacity=1, inner sep=1.5pt},
  ]

  % ================================================================
  % Top: segments shared by A and B
  % ================================================================
  \def\sy{0.0}

  % (4) reduce vertical gap between the "Replica ..." row and segment row
  \node[note, anchor=south west] at (0.05, \sy+0.35) {\textbf{Replica A}};
  \node[note, anchor=south east] at (7.60, \sy+0.35) {\textbf{Replica B}};
  \node[note, anchor=south]      at (3.90, \sy+0.35) {same segmentation in A,B};

  \node[segSmall] (s1) at (0.90, \sy) {1};
  \node[segSmall] (s2) at (1.55, \sy) {2};
  \node[segSmall] (s3) at (2.20, \sy) {3};
  \node[segSmall] (s4) at (2.85, \sy) {4};
  \node[note]     at (3.35, \sy) {$\cdots$};
  \node[segHi, minimum width=1.10cm] (sp) at (4.10, \sy) {seg $p$};
  \node[note]     at (5.15, \sy) {$\cdots$};
  \node[segSmall] (sP) at (6.40, \sy) {$P$};

  % (1) put "same in A,B" together with h_v(p) label
  \draw[->, thick] (sp.south) -- (4.10, -0.95)
    node[arrlbl, pos=0.7, right=3pt, align=center] {$h_v(p)$\\{\scriptsize same in A,B}};

  % ================================================================
  % Section strip with a contiguous W-window
  % ================================================================
  \def\stripX{0.40}
  \def\stripY{-1.80}
  \def\nSec{12}
  \def\cw{0.50}
  \def\stripH{0.36}
  \def\wStart{4}
  \def\Wwin{4}

  \draw[thick] (\stripX,\stripY)
    rectangle ({\stripX+\nSec*\cw},{\stripY+\stripH});
  \foreach \i in {1,...,11} {
    \draw[thin, black!30]
      ({\stripX+\i*\cw},\stripY)
      -- ({\stripX+\i*\cw},{\stripY+\stripH});
  }

  % Blue window fill + border
  \fill[blue!15]
    ({\stripX+\wStart*\cw},\stripY)
    rectangle ({\stripX+(\wStart+\Wwin)*\cw},{\stripY+\stripH});
  \draw[thick, blue!65!black]
    ({\stripX+\wStart*\cw},\stripY)
    rectangle ({\stripX+(\wStart+\Wwin)*\cw},{\stripY+\stripH});

  % W-1 dashed separators inside the window
  \foreach \t in {1,...,3} { % assumes Wwin=4
    \draw[blue!65!black, dashed, line width=0.6pt]
      ({\stripX+(\wStart+\t)*\cw},\stripY)
      -- ({\stripX+(\wStart+\t)*\cw},{\stripY+\stripH});
  }

  % (2) remove the 4 "W" letters inside the window sections
  % (removed)

  % L brace on top
  \draw[decorate, decoration={brace, amplitude=4pt}, black!60]
    ({\stripX},{\stripY+\stripH+0.08})
    -- ({\stripX+\nSec*\cw},{\stripY+\stripH+0.08})
    node[midway, above=4pt, lbl] {$L$ sections};

  % (3) change label text to "W contiguous sections"
  \draw[decorate, decoration={brace, amplitude=4pt, mirror}, blue!65!black]
    ({\stripX+\wStart*\cw},{\stripY-0.08})
    -- ({\stripX+(\wStart+\Wwin)*\cw},{\stripY-0.08})
    node[midway, below=4pt, lbl, text=blue!65!black]
    {$W$ contiguous sections};

  \node[note, anchor=east] at ({\stripX-0.10},{\stripY+\stripH/2})
    {view $v$};

  % (5) shorten the black arrow a bit (end slightly above the brace/grid top)
  \draw[->, thick]
    ({\stripX+(\wStart+\Wwin/2)*\cw}, {\stripY-0.55})
    -- ({\stripX+(\wStart+\Wwin/2)*\cw}, {\stripY-1.07})
    node[arrlbl, midway, right=3pt] {$k$ cells};

  % ================================================================
  % Window-expanded cell plane: W (rows) x m0 (cols)
  % ================================================================
  \def\gx{1.60}
  \def\gy{-2.95}
  \def\cs{0.30}
  \def\nRow{4}  % visual W
  \def\nCol{7}  % visual m0

  % Grid
  \foreach \r in {0,...,3} {
    \foreach \c in {0,...,6} {
      \node[cell, anchor=north west] (c\r\c)
        at ({\gx+\c*\cs},{\gy-\r*\cs}) {};
    }
  }

  % m0 brace (top)
  \draw[decorate, decoration={brace, amplitude=3pt}, black!60]
    ({\gx},{\gy+0.06}) -- ({\gx+\nCol*\cs},{\gy+0.06})
    node[midway, above=3pt, lbl] {$m_0$ cells};

  % W brace (left)
  \draw[decorate, decoration={brace, amplitude=3pt, mirror}, blue!65!black]
    ({\gx-0.10},{\gy})
    -- ({\gx-0.10},{\gy-\nRow*\cs})
    node[midway, left=4pt, lbl, text=blue!65!black] {$W$};

  % Example signed updates
  \fill[red!20]   (c02.north west) rectangle (c02.south east);
  \draw[thick, red!70!black] (c02.north west) rectangle (c02.south east);
  \node[font=\tiny, red!70!black] at (c02) {$+$};

  \fill[blue!20]  (c15.north west) rectangle (c15.south east);
  \draw[thick, blue!70!black] (c15.north west) rectangle (c15.south east);
  \node[font=\tiny, blue!70!black] at (c15) {$-$};

  \fill[black!8]  (c31.north west) rectangle (c31.south east);
  \draw[thick, black!45] (c31.north west) rectangle (c31.south east);
  \node[font=\tiny, black!55] at (c31) {$0$};

  % (6) tighten line spacing on the right text block
  \node[note, anchor=west, align=left]
    at ({\gx+\nCol*\cs+0.25},{\gy-0.10})
    {A-side key: \textcolor{red!70!black}{$+$}};
  \node[note, anchor=west, align=left]
    at ({\gx+\nCol*\cs+0.25},{\gy-0.38})
    {B-side key: \textcolor{blue!70!black}{$-$}};
  \node[note, anchor=west, align=left]
    at ({\gx+\nCol*\cs+0.25},{\gy-0.66})
    {common: cancel $\to 0$};
  \node[note, anchor=west, align=left]
    at ({\gx+\nCol*\cs+0.25},{\gy-0.94})
    {$\Delta_j = A_j - B_j$};

\end{tikzpicture}
\caption{Two-level mapping in one view: segment $p$ is hashed
to $W$ contiguous sections, then each key is hashed to $k$
cells within that window.}
\label{fig:scout-mapping-structure-a}
\end{figure}

%\paragraph{Pipeline}
\noindent\textbf{Pipeline}\quad
Each reconciliation cycle executes four stages; the first
three are fixed-cost, while the fourth adapts to the observed
evidence.

\textbf{(1)~Build.}  Each replica sequentially scans its local key space; every item is inserted into the probe array via the two-level mapping described above.

\textbf{(2)~Exchange \& Diff.}  The two replicas exchange
their probes once; each side subtracts cell-wise to obtain the
difference array
$\Delta_j = C^{(A)}_j - C^{(B)}_j$.
Because both replicas apply identical hash mappings to each
item, any item present on both sides produces equal updates to
the same probe cells.  These contributions cancel under
cell-wise subtraction, so $\Delta$ is driven solely by items
that differ between the two sides: items in $A\setminus B$
contribute $+1$ signed updates, while those in
$B\setminus A$ contribute~$-1$.

\textbf{(3)~Decode.}  For every section~$s$ in each view, the decoder computes the first moment and the raw second moment:
\[
\mathrm{net\_sum}_s = \sum_{j \in s}\Delta_j
,~
\mathrm{sq\_sum}_s = \sum_{j \in s}\Delta_j^2,
\]
then constructs the \textbf{de-drifted energy}:
\begin{equation*}
E_s \;=\; \mathrm{sq\_sum}_s \;-\;
         \frac{(\mathrm{net\_sum}_s)^{2}}{m_0}.
\end{equation*}
$E_s$ is the unnormalized sample variance of the differential
counts within section~$s$.  It serves as the \emph{common signal}
from which all three outputs are derived: triage consumes the
spatial distribution of $\{E_s\}$ (where energy concentrates),
sizing consumes the total mass $\sum_s E_s$ (how much energy
exists), and gating consumes the support dispersion of
$\{E_s\}$ (over how many sections energy spreads).  All three
readings come from the same energy spectrum
without budget
partitioning.  The subtraction of $(\mathrm{net\_sum}_s)^2/m_0$ removes the
first-order sign bias that would otherwise dominate the raw
second moment under side imbalance, so $E_s$ captures
fluctuation energy rather than drift-induced offset (note that
$E_s \ge 0$ by the Cauchy--Schwarz inequality). The full
derivation is given in
Supplement~\ref{sec:appendix-sizing-derivation}.

\textbf{(4)~Decide.} 
Each diagnostic output reads the same section-level energy
spectrum but fuses information across the $V$ views
differently. Sizing computes a per-view scale estimate and averages across views. 
Gating evaluates the concentration metric and occupancy on the energy spectrum (\S\ref{sec:gating}). Triage retains per-view window rankings and intersects the resulting candidate supersets to suppress view-dependent false inclusions.

When all sections in a segment's mapped window carry zero
signal ($E_s{=}0$ and $\mathrm{net\_sum}_s{=}0$) in any view, the segment can be excluded from downstream
verification, because any diff key would have left a
non-zero cell update within the window
(Supplement~\ref{sec:appendix-sizing-zero-energy}).

Sizing emits $\hat{D}$ (\S\ref{sec:dedrift});
$\hat{D}\approx 0$ terminates the round without downstream
verification.  Otherwise, gating tests $K_{\mathrm{eff}}/L$
and occupancy (\S\ref{sec:gating}).
On \texttt{notgate}, triage produces~$\mathcal{C}$ via
sliding-window aggregation and cross-view intersection
(Fig.~\ref{fig:scout-mapping-structure-b}).
On \texttt{gate}, the system suppresses candidate generation
and defers to a conservative fallback---e.g., full-scope exact
verification or a coarse per-range checksum pass over all $P$
segments---whose cost scales with $P$ rather than
$|\mathcal{C}|$.
On \texttt{unknown}, the probe withholds only a shape-based verdict; downstream policy may still invoke capped triage or full-scope fallback
based on $\hat{D}$ and other observables (\S\ref{sec:eval-dedrift}).

\begin{figure}[!htbp]
\centering
% SCOUT probe structure (b): 2 views + reverse mapping + (implicit) intersection + energy peak
\begin{tikzpicture}[
    x=1cm, y=1cm, font=\footnotesize, >=Stealth,
    lbl/.style    = {font=\scriptsize, align=center, inner sep=1.5pt},
    note/.style   = {font=\scriptsize, text=black!60},
    arrlbl/.style = {font=\scriptsize, align=center, fill=white,
                     fill opacity=.92, text opacity=1, inner sep=1.5pt},
    segbox/.style = {draw, rounded corners=1pt, minimum width=0.34cm,
                     minimum height=0.30cm, inner sep=0pt, font=\tiny},
  ]

  % ----------------------------------------------------------------
  % Geometry knobs
  % ----------------------------------------------------------------
  \def\mvX{0.65}
  \def\mvN{12}
  \def\mvW{0.42}
  \def\mvH{0.30}
  \def\mvGap{0.70}
  \pgfmathsetmacro{\stripRight}{\mvX+\mvN*\mvW}

  % ----------------------------------------------------------------
  % Segment row (top)
  % ----------------------------------------------------------------
  \def\segX{\mvX}
  \def\segY{0.45}
  \def\segN{10}
  \def\segW{0.50}

  % segments: 1..10
  % true hot: 5,6 ; halo (normal): 4,7 ; halo (small): 3,8 ; rest: neutral
  \foreach \i in {0,...,9} {
    \pgfmathtruncatemacro{\sid}{\i+1}
    \pgfmathtruncatemacro{\isHot}{(\sid==5 || \sid==6) ? 1 : 0}
    \pgfmathtruncatemacro{\isHalo}{(\sid==4 || \sid==7) ? 1 : 0}
    \pgfmathtruncatemacro{\isHaloSmall}{(\sid==3 || \sid==8) ? 1 : 0}

    \ifnum\isHot=1
      \node[segbox, fill=orange!60, line width=0.8pt]
        (seg\sid) at ({\segX+\i*\segW},{\segY}) {\sid};
    \else
      \ifnum\isHalo=1
        \node[segbox, fill=orange!20, line width=0.6pt]
          (seg\sid) at ({\segX+\i*\segW},{\segY}) {\sid};
      \else
        \ifnum\isHaloSmall=1
          \node[segbox, fill=orange!10, draw=orange!25, line width=0.5pt]
            (seg\sid) at ({\segX+\i*\segW},{\segY}) {\sid};
        \else
          \node[segbox, fill=black!5, draw=black!25]
            (seg\sid) at ({\segX+\i*\segW},{\segY}) {\sid};
        \fi
      \fi
    \fi
  }

  % ----------------------------------------------------------------
  % Two view strips (middle)
  % ----------------------------------------------------------------
  \def\mvYstart{-0.55}

  % View 1: selected window at sections 2–5
  \pgfmathsetmacro{\yyA}{\mvYstart}
  \draw[thick] (\mvX,\yyA) rectangle (\stripRight,{\yyA+\mvH});
  \foreach \i in {1,...,11} {
    \draw[thin, black!25]
      ({\mvX+\i*\mvW},\yyA) -- ({\mvX+\i*\mvW},{\yyA+\mvH});
  }
  % halo-ish neighbors (dashed)
  \fill[orange!10]
    ({\mvX+0*\mvW},\yyA) rectangle ({\mvX+2*\mvW},{\yyA+\mvH});
  \draw[thick, orange!45!black, dashed]
    ({\mvX+0*\mvW},\yyA) rectangle ({\mvX+2*\mvW},{\yyA+\mvH});
  \fill[orange!10]
    ({\mvX+6*\mvW},\yyA) rectangle ({\mvX+7*\mvW},{\yyA+\mvH});
  \draw[thick, orange!45!black, dashed]
    ({\mvX+6*\mvW},\yyA) rectangle ({\mvX+7*\mvW},{\yyA+\mvH});
  % main selected window (solid)
  \fill[orange!35]
    ({\mvX+2*\mvW},\yyA) rectangle ({\mvX+6*\mvW},{\yyA+\mvH});
  \draw[thick, orange!70!black]
    ({\mvX+2*\mvW},\yyA) rectangle ({\mvX+6*\mvW},{\yyA+\mvH});
  \node[note, anchor=east] at ({\mvX-0.06},{\yyA+\mvH/2}) {view 1};

  % View 2: selected window at sections 6–9
  \pgfmathsetmacro{\yyB}{\mvYstart - \mvGap}
  \draw[thick] (\mvX,\yyB) rectangle (\stripRight,{\yyB+\mvH});
  \foreach \i in {1,...,11} {
    \draw[thin, black!25]
      ({\mvX+\i*\mvW},\yyB) -- ({\mvX+\i*\mvW},{\yyB+\mvH});
  }
  % halo-ish neighbors (dashed)
  \fill[orange!10]
    ({\mvX+5*\mvW},\yyB) rectangle ({\mvX+6*\mvW},{\yyB+\mvH});
  \draw[thick, orange!45!black, dashed]
    ({\mvX+5*\mvW},\yyB) rectangle ({\mvX+6*\mvW},{\yyB+\mvH});
  \fill[orange!10]
    ({\mvX+10*\mvW},\yyB) rectangle ({\mvX+10.8*\mvW},{\yyB+\mvH});
  \draw[thick, orange!45!black, dashed]
    ({\mvX+10*\mvW},\yyB) rectangle ({\mvX+10.8*\mvW},{\yyB+\mvH});
  % main selected window (solid)
  \fill[orange!35]
    ({\mvX+6*\mvW},\yyB) rectangle ({\mvX+10*\mvW},{\yyB+\mvH});
  \draw[thick, orange!70!black]
    ({\mvX+6*\mvW},\yyB) rectangle ({\mvX+10*\mvW},{\yyB+\mvH});
  \node[note, anchor=east] at ({\mvX-0.06},{\yyB+\mvH/2}) {view 2};

  % ----------------------------------------------------------------
  % Reverse-mapping arrows (avoid crossing)
  % ----------------------------------------------------------------
  % View1 arrow -> near seg5
  \coordinate (v1c) at ({\mvX+4*\mvW},{\yyA+\mvH});
  \coordinate (s1c) at ({\segX+4*\segW},{\segY-0.18}); % around seg5
  \draw[->, thick, black!55] (v1c) -- (s1c);

  % View2 arrow -> near seg5
  \coordinate (v2c) at ({\mvX+8*\mvW},{\yyB+\mvH});
  \coordinate (s2c) at ({\segX+4*\segW},{\segY-0.18}); % around seg5
  \draw[->, thick, black!55] (v2c) -- (s2c);

  % ----------------------------------------------------------------
  % Intersection marker above segment row: brace over seg5..6 labeled \mathcal{C}
  % (keep the brace, but move \bigcap_v to below seg6)
  \draw[decorate, decoration={brace, amplitude=3.5pt}]
    (seg5.north west) -- (seg6.north east)
    node[midway, above=4pt, lbl] {$\mathcal{C}$};

  % \bigcap_v moved to below seg6
  \node[note, anchor=north] at ($(seg6.south)+(0.25,-0.02)$) {$\bigcap_v$};

  % ----------------------------------------------------------------
  % Energy spectrum (bottom) — simplified labels
  % ----------------------------------------------------------------
  \pgfmathsetmacro{\barY}{\yyB - 1.15}
  \def\barX{\mvX}
  \def\barW{\mvW}
  \def\barScale{0.85}

  % Bars: peak aligned with view-1 selected window (sections 2–5)
  \foreach \idx/\ht in {
    0/0.08, 1/0.10,
    2/0.65, 3/0.80, 4/0.72, 5/0.50,
    6/0.10, 7/0.07, 8/0.09, 9/0.06, 10/0.08, 11/0.07} {
    \pgfmathsetmacro{\bh}{\ht*\barScale}
    \pgfmathtruncatemacro{\isPeak}{(\idx>=2 && \idx<=5) ? 1 : 0}
    \ifnum\isPeak=1
      \fill[orange!45]
        ({\barX+\idx*\barW},\barY)
        rectangle ({\barX+(\idx+1)*\barW},{\barY+\bh});
      \draw[orange!70!black]
        ({\barX+\idx*\barW},\barY)
        rectangle ({\barX+(\idx+1)*\barW},{\barY+\bh});
    \else
      \fill[black!10]
        ({\barX+\idx*\barW},\barY)
        rectangle ({\barX+(\idx+1)*\barW},{\barY+\bh});
      \draw[black!30]
        ({\barX+\idx*\barW},\barY)
        rectangle ({\barX+(\idx+1)*\barW},{\barY+\bh});
    \fi
  }
  \draw[thick] (\barX,\barY) -- ({\barX+12*\barW},\barY);

  % x-label
  \node[note] at ({\barX+6*\barW},{\barY-0.18}) {section $s$};

  % Move E_s to the right of the peak (avoid overlap with left label)
  \node[note, rotate=90, anchor=south]
    at ({\barX+7.35*\barW},{\barY+0.45*\barScale}) {$E_s$};

  % peak brace with only "peak"
  \draw[decorate, decoration={brace, amplitude=3pt}, orange!65!black]
    ({\barX+2*\barW},{\barY+0.80*\barScale+0.08})
    -- ({\barX+6*\barW},{\barY+0.80*\barScale+0.08})
    node[midway, above=2pt, lbl, text=orange!65!black] {peak};

  % Left label: align nicely (same x as strip labels), no overlap
  \node[note, anchor=east, align=left]
    at ({\mvX+0.75},{\barY+0.50})
    {view 1 energy};

\end{tikzpicture}
\caption{From energy peaks to candidate segments:
sliding-window aggregation identifies high-energy regions on
the section axis; reverse mapping and cross-view intersection
yield the candidate set~$\mathcal{C}$.}
\label{fig:scout-mapping-structure-b}
\end{figure}

Beyond the unavoidable single sequential scan over local items,
SCOUT incurs communication proportional to the fixed probe
size and decoding cost of $O(VLm_0)$.
Table~\ref{tab:scout-params} summarizes the core probe
parameters.

\begin{table}[!htbp]
\caption{Quick reference for SCOUT core parameters.
Full symbol table in Supplement~\ref{sec:notations}.}
\label{tab:scout-params}
\centering
\small
\setlength{\tabcolsep}{4pt}
\begin{tabular}{@{}cl@{}}
\toprule
Symbol & Meaning \\
\midrule
$P$ & Total segments ($10^4$--$10^7$) \\
$B$ & Concurrently dirty segments (unknown at runtime) \\
$D$ & Symmetric difference size, $D=D^{+}+D^{-}$ \\
$V$ & Views / independent hash channels (2--4) \\
$L$ & Sections per view (1024--8192) \\
$m_0$ & Cells per section (16--64) \\
$W$ & Window width in contiguous sections (2--4) \\
$k$ & Per-key update redundancy (1--3) \\
\bottomrule
\end{tabular}
\end{table}


\subsection{Observation Geometry for Candidate Compression}
\label{sec:spatial-coupling}

Triage operates on the energy spectrum $\{E_s\}$.
Its ability to compress the candidate set depends on whether
dirty-segment energy remains localized on the section axis.
Localized energy produces separable peaks and enables
candidate compression; dispersed energy blends into
background noise and reduces triage to broad ranking.


\paragraph{Contiguous-interval mapping and sliding-window aggregation}
To preserve locality, SCOUT maps each segment to $W$
\textbf{contiguous} sections on the section axis.  Triage
therefore scores \emph{windows} rather than individual
sections.  For each starting position~$u$, SCOUT computes
\[
\Omega_u = \sum_{s=u}^{u+W-1} E_s,
\]
the total energy in the width-$W$ window beginning at~$u$.
This sliding-window aggregation pools the energy that a single
segment deposits across its $W$ adjacent sections, strengthening
localized peaks while suppressing unstructured background
noise---a width-$W$ matched filter in signal-processing terms.

SCOUT then ranks all windows by~$\Omega_u$ and keeps the top
$
t = \left\lceil r_{\mathrm{top}} \cdot (L-W+1) \right\rceil
$
windows, where $r_{\mathrm{top}} \in (0,1]$ controls the
candidate budget. Each selected window is then reverse-mapped
to the segments whose supports overlap it
(Fig.~\ref{fig:scout-mapping-structure-b}), producing a
view-local candidate superset.
The final candidate set~$\mathcal{C}$ is obtained by
intersecting these supersets across all $V$ views, which
suppresses false inclusions exponentially under approximately
independent views while preserving true positives with high
probability.

The compression effect above depends on \emph{contiguity}: if
the $W$ target sections were scattered randomly instead of
placed contiguously, the sliding window would lose its locality
advantage and the scheme would reduce to independent
per-section ranking (\S\ref{sec:eval-spatial}).

\paragraph{Scaling laws for candidate-set size}
Under the uniform-hashing model (\S\ref{sec:overview}),
if concentrated differences produce separable signals,
the expected candidate-set size after intersecting across $V$
views is:
\begin{equation}\label{eq:cand-expectation}
\mathbb{E}[|\mathcal{C}|]
  \;\approx\; B \;+\; (P-B)\cdot p_{\mathrm{halo}}^{\,V},
\end{equation}
where $B$ is the number of truly dirty segments and the second
term captures the \emph{halo} contribution: clean segments whose
windows geometrically overlap with a selected dirty window are
erroneously included in the candidate set.
A clean segment's false-inclusion probability in a single view
is
\begin{equation}\label{eq:halo-prob}
p_{\mathrm{halo}}
  \;\approx\; \frac{t(2W-1)}{L}
  \quad(\text{when } t(2W-1) \ll L),
\end{equation}
where $2W{-}1$ is the overlap neighborhood width of each
selected window on the section axis.
This framework implies three scaling laws:
(i)~cross-view intersection raises $p_{\mathrm{halo}}$ to the
$V$-th power, making $V$ the most effective lever against
candidate inflation;
(ii) the halo term grows with the number of clean segments,
while larger $L$ reduces per-view overlap and suppresses false
inclusion;
(iii)~$W$ improves signal aggregation but simultaneously
inflates the halo via the $2W{-}1$ factor, so there is no
uniformly optimal~$W$
(derivations in Supplement~\ref{sec:appendix-candidate-expectation}
and~\ref{sec:appendix-w-pricing}).
All three laws require that dirty windows remain separable
on the section axis; when $L/B$ is too small, their supports
overlap and candidate cost degrades sharply.
This operating boundary is characterized empirically in
\S\ref{sec:eval-operating} and modeled in
Supplement~\ref{sec:appendix-operating-theory}.


\subsection{Drift-Stable Sizing from the Shared Substrate}
\label{sec:dedrift}

The de-drifted energy~$E_s$ (\S\ref{sec:overview}) removes
the first-order sign bias that would otherwise cause sizing
to overestimate by orders of magnitude under drift; we now
derive the estimator built on it.
Let $M = L \cdot m_0$ denote the total cell count
of a single view.
Under the same uniform-hashing model, the expected total
de-drifted energy is:
\[
\mathbb{E}\!\left[\sum_s E_s\right]
  \;=\; D \cdot k\!\left(1 - \frac{k}{M}\right).
\]
Mean subtraction cancels the drift factor~$(2q{-}1)^2$
exactly: both the squared-mean component of~$\mathrm{sq\_sum}$
and the correction term~$\mathrm{net\_sum}^{2}/m_0$ scale
quadratically with the mean sign bias~$(2q{-}1)$, and their
coefficients match, so after subtraction the expectation does
not depend on~$q$.
The sizing estimator is
\[
  \hat{D} \;=\; \frac{\sum_s E_s}{\,k\bigl(1 - k/M\bigr)\,},
\]
where the denominator $1{-}k/M$ corrects for energy dilution
from cross-key collisions within the same cell.

Two minor bias sources remain.
First, $E_s$ as a sample variance exhibits finite-sample
shrinkage, removable by a Bessel correction factor
$m_0/(m_0{-}1)$.
Second, when differences are highly concentrated in a few
sections the global uniformity approximation introduces a
systematic underestimate of roughly $1$--$2\%$.
After both corrections the measured relative bias converges to
within approximately $3\%$ across all tested configurations (\S\ref{sec:eval-dedrift}).

Multi-view fusion provides additional robustness:
averaging across approximately independent views drives
the standard deviation down as $1/\sqrt{V}$, making $V$
and $m_0$ the two main tuning knobs for sizing precision.
Derivations covering drift cancellation, bias
corrections, and multi-view fusion are given in
Supplement~\ref{sec:appendix-sizing-derivation}
and~\ref{sec:appendix-sizing-bias}.


\subsection{Admission and Abstention: When Not to Triage}
\label{sec:gating}

The gating decision rests on two observable quantities
derived from the energy spectrum: how broadly the energy is
spread, and whether the sketch carries enough statistical
support to make that measurement reliable.

\paragraph{Concentration metric}
We measure energy dispersion by the inverse Herfindahl index
\begin{equation*}
\frac{K_{\mathrm{eff}}}{L}
  \;=\; \frac{1}{L}\,\frac{\bigl(\sum_s E_s\bigr)^{2}}
        {\sum_s E_s^{2}}
  \;\in\; [1/L,\,1].
\end{equation*}
$K_{\mathrm{eff}}/L$ measures how many sections carry appreciable energy. It is invariant to overall discrepancy scale and to drift when computed from de-drifted energy. Under concentrated
faults, energy resides in roughly $BW$ sections, placing $K_{\mathrm{eff}}/L$ at the scale of $BW/L$, typically far below~$1$; under diffuse differences with sufficient occupancy it
approaches~$1$. In the concentrated regime, $K_{\mathrm{eff}}/L$ also governs candidate inflation: a larger value requires a more generous truncation budget to preserve recall, widening
halo overlap so that candidate cost grows as $p_{\mathrm{halo}}^V$. $K_{\mathrm{eff}}/L$ therefore serves not only as a shape statistic but as a cost-oriented proxy for whether triage
remains economical. Because concentrated and diffuse discrepancies are separated by roughly an order of magnitude in this metric, $\tau_{\text{gating}}$ needs only to fall between the two
regimes rather than be tuned precisely (derivation in Supplement~\ref{sec:appendix-gating-separation}; empirical sensitivity in \S\ref{sec:eval-dedrift}).

\paragraph{Observability guard-rail and three-state output}
The scale separation above requires sufficient statistical
support in the energy spectrum. When occupancy is low,
diffuse cases can collapse to support sizes comparable to
concentrated ones, making $K_{\mathrm{eff}}/L$ unreliable.
We define occupancy as the cross-view fraction of non-empty
sections:
\begin{equation}\label{eq:occupancy}
\mathrm{occ}
  \;=\; \frac{1}{V}\sum_{v=1}^{V}\;
        \frac{1}{L}\sum_{s=1}^{L}
        \mathbf{1}\!\bigl[E_s^{(v)} > 0\bigr],
\end{equation}
where $E_s^{(v)}$ denotes $E_s$ evaluated on view~$v$.
SCOUT uses a three-state gating rule:
$g=\texttt{unknown}$ if $\mathrm{occ}<\tau_{\mathrm{occ}}$;
otherwise, $g=\texttt{notgate}$ if $K_{\mathrm{eff}}/L\le\tau_{\mathrm{gating}}$,
and $g=\texttt{gate}$ otherwise.
Neither threshold requires narrow tuning.
$\tau_{\text{gating}}$ only needs to fall between the
concentrated and diffuse regimes of $K_{\mathrm{eff}}/L$,
a gap that remains wide whenever occupancy is sufficient
(\S\ref{sec:eval-dedrift}).
$\tau_{\text{occ}}$ only needs to exceed the typical
concentrated-fault occupancy~$BW/L$ by a moderate margin
to preserve scale separation within the observable region.
